\documentclass{article}

% if you need to pass options to natbib, use, e.g.:
%     \PassOptionsToPackage{numbers, compress}{natbib}
% before loading neurips_2025

\usepackage[main, final]{neurips_2025}

\usepackage[utf8]{inputenc} % allow utf-8 input
\usepackage[T1]{fontenc}    % use 8-bit T1 fonts
\usepackage{hyperref}       % hyperlinks
\usepackage{url}            % simple URL typesetting
\usepackage{booktabs}       % professional-quality tables
\usepackage{amsfonts}       % blackboard math symbols
\usepackage{nicefrac}       % compact symbols for 1/2, etc.
\usepackage{microtype}      % microtypography
\usepackage{xcolor}         % colors
\usepackage{tabularx}
\usepackage{float}

\makeatletter
\renewcommand{\@noticestring}{}
\makeatother

\title{Solving Enhanced Frozen Lake via Memory-Based and Distributional Reinforcement Learning}


\author{%
  Shengkai (Kevin) Tao, Xiyuan (Sophia) Shen \\
  Department of Data Science \\
  UNC Chapel Hill \\
  \texttt{kevintao@unc.edu} \\
}

\begin{document}

\maketitle

\begin{abstract}
While standard RL models excel in static environments, real-world applications demand resilience against uncertainty. Whether it is a search-and-rescue drone or autonomous underwater devices, agents must handle partial observability and dynamic environmental shifts. Our project is motivated by the need for agents that don't just solve a map but understand its risks. To solve the common issue where high penalties stifle learning, we integrate memory-based architectures with exploration driven by curiosity. We use an Enhanced Frozen Lake simulation to bridge this gap, testing whether an agent can remain adventurous enough to find the goal while remaining smart enough to survive the challenges.
\end{abstract}

\section{Problem Statement \& Motivation}

\textbf{Motivation:} While standard RL models excel in static environments, real-world applications demand resilience against uncertainty. Whether it is a search-and-rescue drone or autonomous underwater devices, agents must handle partial observability and dynamic environmental shifts. Our project is motivated by the need for agents that don't just solve a map but understand its risks. To solve the common issue where high penalties stifle learning, we integrate memory-based architectures with exploration driven by curiosity. We use an Enhanced Frozen Lake simulation to bridge this gap, testing whether an agent can remain adventurous enough to find the goal while remaining smart enough to survive the challenges.

\textbf{Frozen Lake:} Agent travelling on a frozen lake (slippery) grid. It starts on a cell and needs to reach the goal. There are holes in the grid that the agent needs to avoid. Can only move up, down, left, right.

\textbf{Enhanced Frozen Lake:} Our project adds 3 complications on top of the standard Frozen Lake. This now makes the game a Partially Observable Markov Decision Process (POMDP). Each complication is addressed with a specific method, which will be tested and compared against.

\begin{table}[H]
    \centering
    \caption{Challenges and Proposed Solutions}
    \begin{tabularx}{\textwidth}{l X l}
        \toprule
        \textbf{Challenge} & \textbf{Why It Is Hard} & \textbf{Proposed Solution} \\
        \midrule
        Partial Observability & Agent only sees a 5x5 window: it cannot use a table keyed on full map state. & DRQN (LSTM memory) \\
        Non-Stationary Dynamics & Introducing a "wind" variable that changes every 1000 episodes: transition probabilities shift; a static policy will degrade. & QR-DQN (distributional RL) \\
        Costly Exploration & Agent navigates a randomly generated map. Without curiosity it learns to play safe and stops exploring. & RND-DQN (exploration) \\
        \bottomrule
    \end{tabularx}
\end{table}

\section{Background and Related Work}

\subsection{Deep Q-Network (DQN) Baseline}
``Human-Level Control through Deep Reinforcement Learning'' \cite{mnih2015human}

This paper establishes the baseline agent we will compare all methods against. DQN is an upgrade to Q-learning in that instead of using a lookup table that only stores 1 Q value for each state, it combines convolutional neural network to output 1 Q value per action per state. The training will build an intuition that can be carried to future, similar states, making it possible to learn in large state spaces.

This paper's proposed DQN method is a breakthrough on previous RL methods and achieves to play 49 Atari games at human level. This training is possible because of 1) experience relay that trains on randomly sampled mini batches of transitions stored in a buffer, preventing the network from overfitting to recent consecutive experiences and stabilizing learning across diverse situations 2) a target network that provides a fixed reference point for computing training targets, preventing the feedback loop that occurs when the network chases a constantly shifting goal, which would otherwise cause training to diverge.

\textbf{Limitations:} DQN struggles with partial observability, uncertainty in stochastic environments (random wind variable), adapting policies if environment dynamic changes (change of wind direction).

\subsection{Deep Recurrent Q-Network (DRQN)}
``Deep Recurrent Q-Learning for Partially Observable MDPs'' \cite{hausknecht2017}

This research proposed the DRQN that addresses the shortcomings of requiring full observability in DQN. DRQN replaces the first fully connected layer in DQN with a Long Short-Term Memory (LSTM) recurrent layer that enables the agent to make decisions based on what it sees now and everything it remembers seeing (memory). LSTM works by maintaining a hidden state that gets passed from one timestep to the next and updated at each step. It can decide what to forget (Forget Gate), what to remember (Input Gate), and what to use right now (Output Gate).

\textbf{Limitations:} Slow to train because LSTM is sequential and may not outperform DQN on very simple POMDPs.

\subsection{Distributional RL with Quantile Regression (QR-DQN)}
``Implicit Quantile Networks for Distributional Reinforcement Learning'' \cite{dabney2018}

\subsection{RL in Non-Stationary Environments}

\subsection{RND}
``Exploration by Random Network Distillation'' \cite{burda2018}

\section{Methods}

\textbf{Baseline: Q-learning and DQN} \\
Maps: We check feasibility of randomly generated maps with BFS. 

\textbf{DRQN (Deep Recurrent Q-Network)} is the foundation. The core idea is that the agent only sees a 5$\times$5 window around itself, so it can't see the whole map --- this is called a Partially Observable environment. A regular DQN would be helpless here because each observation alone doesn't tell you enough. DRQN solves this by adding an LSTM (a type of memory unit) so the agent remembers what it saw in previous steps. As it walks around, it mentally builds up a picture of the map from its limited field of view. This is your baseline method --- the one everything else is compared against.

``While Mnih et al. employ a convolutional network suited for pixel-based inputs, our implementation replaces this with tile embeddings appropriate for symbolic grid observations.''

The original DRQN only takes in visual observations. Your implementation adds wind direction as a second explicit input through a separate embedding layer. That's an enhancement not present in the original.

Your code directly implements this, storing full episodes in the EpisodeReplayBuffer. The key difference is that your buffer stores episodes rather than individual transitions, because your LSTM needs sequential traces to build its hidden state. Your code does this too --- the target\_net is copied from policy\_net every few episodes rather than updated continuously.

\textbf{QR-DRQN (Quantile Regression DRQN)} takes DRQN and makes it smarter about uncertainty. Standard DRQN asks ``what's the expected reward from this action?'' --- a single number. But in your environment, the wind randomly pushes the agent in unexpected directions, so outcomes are highly variable. QR-DRQN instead asks ``what's the full distribution of possible rewards?'' --- it predicts 51 different percentile estimates (like: in the worst 10\% of cases I get X, in the median case I get Y, in the best 10\% of cases I get Z). This lets it make more cautious decisions when the environment is unpredictable. The theoretical justification comes from the \cite{dabney2018} paper in your background document.

\textbf{RND-DRQN (Random Network Distillation + DRQN)} takes DRQN and adds a curiosity bonus. The problem it addresses is exploration: on a 8×8 grid with a -0.1 living penalty, the agent gets punished for every step it takes, so it might give up exploring and just stand still or repeat safe loops. RND fixes this by giving the agent a bonus reward for visiting states it hasn't seen before. It works by having two small networks --- one frozen random network and one learnable one. When the learnable network can't predict what the frozen one does, that means ``novel state $\rightarrow$ give curiosity reward.'' Over time, as the agent visits everywhere, the curiosity reward naturally fades away. This is based on the \cite{burda2018} paper in your background doc.

So the purpose of having all three is scientific comparison: does memory alone (DRQN) solve the problem? Does accounting for uncertainty (QR-DRQN) do better under wind? Does curiosity-driven exploration (RND-DRQN) help the agent find the goal faster? The evaluate\_models.py script answers exactly this.

\subsection{What ``Wind as Explicit Second Input'' Means}
In the original v1 code, the agent only received the 3×3 grid observation. The wind variable existed in the environment, but the agent had no way of knowing what the wind was doing --- it could only feel its effects indirectly when it got pushed in unexpected directions.

In map8\_v5, your teammate changed this so the current wind direction is passed directly into the model as a second input alongside the visual observation. In the code, there's a separate wind\_embedding layer --- the wind direction (0, 1, 2, or 3) gets converted into a 16-dimensional vector, which is then concatenated with the visual features before being fed into the LSTM.

The effect is significant: now when the wind shifts, the agent isn't surprised --- it ``knows'' and can factor that into its decisions. For example, if wind is pushing left, a smart agent should compensate by taking more right-leaning paths. Without this explicit input, the agent has to infer wind direction purely from experience, which is much harder.

This is directly relevant to the non-stationary dynamics challenge your proposal described. The supporting literature is the Padakandla et al. (2019) paper in your background doc on RL in non-stationary environments.
\section{Experiments and Results}

\subsection{Experimental Setup}
The experiments were conducted in the \textbf{Enhanced Frozen Lake} environment on an $8 \times 8$ grid. Unlike the standard environment, the agent is restricted to a $5 \times 5$ partially observable window. We introduced two stochastic elements: a dynamic "wind" variable that shifts every 1,000 episodes and a living penalty of $-0.1$ per step to discourage aimless wandering. The agent receives a reward of $+30$ for reaching the goal and $-10$ for falling into a hole.

\subsection{Baselines: Tabular Q-Learning and DQN}
We first established performance baselines using Tabular Q-Learning and a Standard Multi-Input Deep Q-Network (DQN). As expected, Tabular Q-Learning struggled significantly, achieving a success rate of only $17.0\%$. This failure is attributed to the high-dimensional state space created by partial observability; the number of possible $5 \times 5$ window configurations combined with wind directions exceeds what a lookup table can effectively map.

Standard DQN improved performance to $63.0\%$ by utilizing tile embeddings and concatenating wind information. However, its lack of memory remains a critical bottleneck. Without recurrency, the agent cannot distinguish between identical local observations that require different actions based on their position relative to the global map.

\subsection{The Role of Memory and Wind Awareness: DRQN}
By introducing an LSTM layer in the Deep Recurrent Q-Network (DRQN), the success rate jumped to $94.0\%$. The hidden state allowed the agent to integrate sequences of observations, effectively "mapping" its surroundings as it moved. Furthermore, we found that providing the \textbf{wind direction as an explicit second input} to the model had a significant effect on performance. By concatenating a learned wind embedding with the visual features, the agent could anticipate environmental shifts rather than inferring them purely through trial and error. This allowed the DRQN to proactively compensate for the wind's influence, leading to a much more stable and reliable policy compared to the DQN baseline. This confirms that both memory and explicit awareness of non-stationary dynamics are essential for solving the Enhanced Frozen Lake.

\subsection{Advanced Hybrid Models: QR-DRQN and RND-DRQN}
While DRQN performed well, it occasionally failed when the wind shifted or when it got stuck in local reward minima. To address these, we implemented two hybrid models:
\begin{itemize}
    \item \textbf{QR-DRQN}: By predicting a distribution of returns via Quantile Regression, the agent became more robust to the aleatoric uncertainty introduced by the wind. It achieved a success rate of $97.0\%$.
    \item \textbf{RND-DRQN}: By adding a curiosity-driven exploration bonus, the agent explored the map more thoroughly during early training, reaching the goal faster and more reliably. It also achieved a success rate of $97.0\%$, with slightly higher average rewards due to more efficient pathfinding.
\end{itemize}

\subsection{Evaluation Methodology}
To ensure a fair and scientifically rigorous comparison between the models, we implemented a custom evaluation script, \texttt{evaluate\_models.py}. The core logic of our evaluation involves testing each agent over 100 episodes using \textbf{fixed seeds} for map generation. By fixing the global seed (incrementing from 2000), we ensure that all agents—from the tabular baseline to the advanced hybrid models—are tested on the exact same sequence of map configurations and starting positions.

We also calculated an \textbf{Oracle Baseline} using a Breadth-First Search (BFS) algorithm. This Oracle has perfect visibility of the full map and calculates the optimal shortest path from start to goal. The Oracle's reward is calculated by applying the same living penalties and success bonuses as the learning agents, providing a theoretical upper bound for performance in each specific map instance.

Each model is evaluated on three primary metrics:
\begin{itemize}
    \item \textbf{Success Rate}: The percentage of episodes where the agent successfully reached the goal.
    \item \textbf{Average Reward}: The total reward accumulated per episode, averaged across all test runs.
    \item \textbf{Average Steps to Success}: The mean number of steps required to reach the goal in successful episodes, serving as a proxy for path efficiency.
\end{itemize}

\subsection{Comparison of Results}
The results of our evaluation over 100 test episodes are summarized in Table \ref{results-table}.

\begin{table}[H]
    \centering
    \caption{Performance Comparison of RL Agents on Enhanced Frozen Lake}
    \label{results-table}
    \begin{tabularx}{\textwidth}{l X c c}
        \toprule
        \textbf{Model} & \textbf{Key Architecture} & \textbf{Success Rate} & \textbf{Avg Reward} \\
        \midrule
        Oracle (BFS) & Shortest Path (Full Map) & $100.0\%$ & $30.70$ \\
        \midrule
        Tabular Q-Learning & State-Wind Lookup Table & $17.0\%$ & $-2.53$ \\
        Standard DQN & Tile Embeddings + MLP & $63.0\%$ & $8.65$ \\
        Standard DRQN & LSTM Memory & $94.0\%$ & $19.24$ \\
        \textbf{QR-DRQN} & \textbf{Distributional + LSTM} & $\mathbf{97.0\%}$ & $\mathbf{19.81}$ \\
        \textbf{RND-DRQN} & \textbf{Curiosity + LSTM} & $\mathbf{97.0\%}$ & $19.76$ \\
        \bottomrule
    \end{tabularx}
\end{table}

\section{Reproducibility and code quality or proof writing}
...
\section{Discussion, Limitation, and Future Directions}
\section{Conclusion}

\bibliographystyle{plain}
\bibliography{Project}

\end{document}
